\section{Introduction}
\label{sec:introduction}
% Fig. 1: overview of the evaluated pipeline (drawn in TikZ; no data in this figure).
\begin{figure*}[t]
\centering
\begin{tikzpicture}[
    font=\footnotesize,
    >={Stealth[length=5pt]},
    box/.style={draw=inkgray, rounded corners=2pt, align=center, inner sep=3pt, minimum height=1.25cm, fill=white},
    lg/.style={box, draw=lgblue, fill=lgblue!7, line width=0.6pt},
    fail/.style={box, draw=failred, fill=failred!6},
    note/.style={font=\scriptsize, text=inkgray, align=center},
    flow/.style={->, draw=inkgray, line width=0.6pt},
    qflow/.style={->, draw=lgblue, line width=0.6pt, dashed},
    node distance=0.36cm,
  ]
  \node[box, text width=1.7cm] (ecg) {\textbf{12-lead ECG}\\ $12\times1000$ samples\\ 10\,s at 100\,Hz};
  \node[fail, right=of ecg, text width=2.0cm] (fail) {\textbf{Lead failure}\\ (simulated)\\ unplug, flatline,\\ EMG, wander, clip};
  \node[box, right=of fail, text width=2.2cm] (enc) {\textbf{Per-lead encoder}\\ shared ResNet-1D\\ + lead embedding\\ $\mathbf{h}_j\in\mathbb{R}^{128}$};
  \node[lg, right=of enc, text width=2.45cm] (mala) {\textbf{Lead attention}\\ 2 layers, 4 heads\\ key bias $\log q_j$\\ masked if $q_j<0.02$};
  \node[lg, right=of mala, text width=1.65cm] (pool) {\textbf{Pooling}\\ weighted by $q_j$\\ $\rightarrow\ \mathbf{z}$};
  \node[lg, right=of pool, text width=2.0cm] (head) {\textbf{Evidential head}\\ Beta$(\alpha_c,\beta_c)$\\ per class $\rightarrow$\\ $p_c$ and $U$};
  \node[box, right=of head, text width=1.75cm, draw=failred, line width=0.6pt] (abst) {\textbf{Abstain} if\\ $U>\tau$ or $\sum_j q_j<3$\\ else 5 labels};

  \node[lg, above=0.6cm of mala, text width=5.2cm, minimum height=0.75cm] (sqa)
       {\textbf{Signal quality assessor} (frozen per-lead CNN):\\ reliability $q_j\in[0,1]$ for every lead};

  \draw[flow] (ecg) -- (fail);
  \draw[flow] (fail) -- (enc);
  \draw[flow] (enc) -- (mala);
  \draw[flow] (mala) -- (pool);
  \draw[flow] (pool) -- (head);
  \draw[flow] (head) -- (abst);
  \draw[flow] (fail.north) |- (sqa.west);
  \draw[qflow] (sqa.south) -- (mala.north);
  \draw[qflow] (sqa.east) -| (pool.north);
  \draw[qflow] (sqa.east) -| (abst.north);

  \node[note, below=0.3cm of mala, text width=5.4cm] (abl)
       {\textbf{Ablation:} un-gated ($q\equiv1$), hard ($\mathbf{1}[q\ge0.5]$) or soft ($q$) gating;
        sigmoid or evidential head};
  \draw[densely dotted, draw=inkgray] (abl.north) -- (mala.south);
  \node[note, below=0.3cm of fail, text width=4.2cm, xshift=1.2cm]
       {\textbf{Baselines} receive the same corrupted ECG (removed leads set to zero) and ignore $q$};
\end{tikzpicture}
\caption{The evaluated pipeline. Lead failure is simulated in training and at test time. A frozen quality assessor
scores every lead; its score $q_j$ biases the lead-attention keys, weights the pooling and enters one of the two
abstention rules. The ablations (bottom) remove the gating or replace the evidential head with a sigmoid head. The
baselines receive the same corrupted ECG, with removed leads set to zero.}
\label{fig:overview}
\end{figure*}


\IEEEPARstart{T}{he} standard electrocardiogram (ECG) records ten seconds from ten electrodes and derives twelve
leads, and its interpretation assumes that each electrode sits where the recording conventions place it
\cite{kligfield2007recommendations}. In emergency departments, ambulances and intensive-care units this assumption
does not always hold. Electrodes lose contact, cables are pulled during transfers, muscle activity adds noise, and
cables can be connected to the wrong electrodes, which can imitate disease \cite{batchvarov2007incorrect}. An
automated classifier used in these settings therefore receives recordings that differ from those it was trained
on.

Deep networks trained on large collections of 12-lead ECGs classify common abnormalities with high accuracy
\cite{ribeiro2020automatic}; on PTB-XL the reference benchmark reports a macro-averaged area under the receiver
operating characteristic curve (macro-AUROC) of about 0.93 for five diagnostic superclasses
\cite{strodthoff2021benchmark}. A convolutional network has no explicit representation of a missing or corrupted
channel, so it returns a diagnosis whatever the state of the input. Three kinds of remedy have been used. The
training data can include recordings with leads removed, as in random lead masking \cite{oh2022leadagnostic},
modality dropout in multimodal learning \cite{neverova2016moddrop}, and the reduced-lead models of the
PhysioNet/Computing in Cardiology Challenge 2021 \cite{reyna2021willtwodo,nejedly2022ensemble}. Information about
signal quality, long studied for ECG acceptability and alarm reduction \cite{clifford2012sqi,behar2013sqi}, can be
passed to the model so that it relies less on poor leads. And the model can estimate its own uncertainty, for
example with evidential deep learning \cite{sensoy2018edl}, so that the system can abstain.

These remedies are often combined. If a quality-aware or uncertainty-aware model is trained on corrupted recordings
and compared with a baseline trained on clean ones, any difference mixes the effect of the architecture with that
of the training data. We therefore ask a narrower question: once a classifier is trained with simulated lead
failure, does gating its lead attention by a per-lead quality estimate, or adding an evidential output that can
abstain, improve robustness or make its uncertainty more useful?

We built a pipeline that contains both components (Fig.~\ref{fig:overview}). A frozen signal-quality assessor
scores every lead, a lead-attention network uses the scores to down-weight or exclude unreliable leads, and a
per-class Beta-evidential output reports an uncertainty that drives abstention. We trained this model, three
ablations that remove the gating or replace the evidential output, and five baselines under one training budget,
and tested all of them on the same deterministically corrupted copies of the PTB-XL test fold, with an oracle
quality control and paired bootstrap intervals. The gated and evidential models were trained only with
corruption, so the study measures what these components add to corruption training, not what they would achieve
without it.

The contributions are three empirical results on PTB-XL with simulated lead failure:
\begin{itemize}
  \item Training with simulated lead failure accounts for most of the difference between models under lead loss:
        the same residual network (ResNet) gains \gainDropKTen{} macro-AUROC with ten of twelve leads removed and
        loses nothing on complete recordings (Section~\ref{sec:res-dropout}).
  \item With that training, hard and soft quality gating show no measurable advantage over un-gated lead
        attention, and the full system is not significantly better than a corruption-trained ResNet in any of
        \bootNRows{} paired comparisons; it is worse under baseline wander and strong muscle noise
        (Section~\ref{sec:res-gating}).
  \item The evidential uncertainty ranks errors no better than the model's predicted probabilities and does not
        increase as leads are removed, while discrimination falls; abstention beyond the clean-data rate comes only
        from a rule that counts usable leads (Section~\ref{sec:res-edl}).
\end{itemize}
The code, trained models and result files are publicly available (see the first-page footnote).
